%% LyX 2.0.3 created this file.  For more info, see http://www.lyx.org/.
%% Do not edit unless you really know what you are doing.
\documentclass[10pt,spanish]{article}
\usepackage[utf8x]{inputenc}
\usepackage[letterpaper]{geometry}
\geometry{verbose}
\usepackage{amsmath}
\usepackage{amssymb}
\usepackage{graphicx}

\makeatletter

%%%%%%%%%%%%%%%%%%%%%%%%%%%%%% LyX specific LaTeX commands.
%% Because html converters don't know tabularnewline
\providecommand{\tabularnewline}{\\}

%%%%%%%%%%%%%%%%%%%%%%%%%%%%%% User specified LaTeX commands.

\usepackage{ucs}\usepackage[spanish]{babel}
\usepackage{amsfonts}\usepackage{colortbl}% = = = = = = = = = = = = = = = = = = = = = = = = =
% INICIO EL DOCUMENTO
% = = = = = = = = = = = = = = = = = = = = = = = = =



\usepackage{babel}
\addto\shorthandsspanish{\spanishdeactivate{~<>}}





\usepackage{babel}
\addto\shorthandsspanish{\spanishdeactivate{~<>}}



\usepackage{babel}
\addto\shorthandsspanish{\spanishdeactivate{~<>}}





\usepackage{babel}
\addto\shorthandsspanish{\spanishdeactivate{~<>}}





\usepackage{babel}
\addto\shorthandsspanish{\spanishdeactivate{~<>}}





\usepackage{babel}
\addto\shorthandsspanish{\spanishdeactivate{~<>}}





\usepackage{babel}
\addto\shorthandsspanish{\spanishdeactivate{~<>}}





\usepackage{babel}
\addto\shorthandsspanish{\spanishdeactivate{~<>}}





\usepackage{babel}
\addto\shorthandsspanish{\spanishdeactivate{~<>}}



\usepackage{babel}
\addto\shorthandsspanish{\spanishdeactivate{~<>}}





\usepackage{babel}
\addto\shorthandsspanish{\spanishdeactivate{~<>}}





\usepackage{babel}
\addto\shorthandsspanish{\spanishdeactivate{~<>}}





\usepackage{babel}
\addto\shorthandsspanish{\spanishdeactivate{~<>}}





\usepackage{babel}
\addto\shorthandsspanish{\spanishdeactivate{~<>}}





\usepackage{babel}
\addto\shorthandsspanish{\spanishdeactivate{~<>}}





\usepackage{babel}
\addto\shorthandsspanish{\spanishdeactivate{~<>}}



\usepackage{babel}
\addto\shorthandsspanish{\spanishdeactivate{~<>}}

\makeatother

\usepackage{babel}
\addto\shorthandsspanish{\spanishdeactivate{~<>}}

\begin{document}

% = = = = = = = = = = = = = = = = = = = = = = = = =
%    INICIO DE TRAYECTORIA
% = = = = = = = = = = = = = = = = = = = = = = = = =
\subsection{CUC2.0: Gestionar Proveedores}

\textbf{\large Resumen}\\
 {\large {} {} Este caso de uso tiene como objetivo permitir al
usuario de compras visualizar los proveedores que existen en el sistema, y le muestra opciones para registrar, modificar, eliminar, desbloquear y filtrar los resultados.}\\


% = = = = = = = = = = = = = = = = = = = = = = = = =
%    INICIO DE TABLA
% = = = = = = = = = = = = = = = = = = = = = = = = = 
\begin{table}[!ht]
\begin{centering}
\begin{tabular}{|c||c|l|}
\hline 
\multicolumn{2}{|c|}{Caso de Uso:} & CUC2.0: Gestionar Proveedores\tabularnewline
\hline 
\multicolumn{3}{|>{\columncolor[gray]{0.7}}c}{Resumen de Atributos}\tabularnewline
\hline 
\multicolumn{2}{|c|}{Autor} & \multicolumn{1}{p{10cm}||}{Jiménez López Carlos Adrian}\tabularnewline
\hline 
\multicolumn{2}{|c|}{Actor} & \multicolumn{1}{p{10cm}||}{Usuario de Compras}\tabularnewline
\hline 
\multicolumn{2}{|c|}{Propósito} & \multicolumn{1}{p{10cm}||}{Gestionar la información de los proveedores que hay en el sistema}\tabularnewline
\hline 
\multicolumn{2}{|c|}{Entradas} & \multicolumn{1}{p{10cm}||}{Ninguna}\tabularnewline
\hline 
\multicolumn{2}{|c|}{Salidas} & \multicolumn{1}{p{10cm}||}{Pantalla PC1}\tabularnewline
\hline 
\multicolumn{2}{|c|}{Pre-condiciones} & \multicolumn{1}{p{10cm}||}{Iniciar sesión como usuario de compras}\tabularnewline
\hline 
\multicolumn{2}{|c|}{Pos-condiciones} & \multicolumn{1}{p{10cm}||}{Ninguna}\tabularnewline
\hline 
\multicolumn{2}{|c|}{Errores} & \multicolumn{1}{p{10cm}||}{MSG4 Error al conectar con la base de datos.}\tabularnewline
\hline 
\multicolumn{2}{|c|}{Tipo} & \multicolumn{1}{p{10cm}||}{Primario}\tabularnewline
\hline 
\multicolumn{2}{|c|}{Fuente} & \multicolumn{1}{p{10cm}||}{}\tabularnewline
\hline 
\end{tabular}
\par\end{centering}

\caption{Caso de Uso Compras 2.0: Gestionar Proveedores}


\label{CUC1.0} 
\end{table}


% = = = = = = = = = = = = = = = = = = = = = = = = =
%	FIN DE TABLA
% = = = = = = = = = = = = = = = = = = = = = = = = =


\textbf{\large Trayectorias del CU}{\large \par}

\textit{\large Trayectoria Principal}{\large {} }{\large \par}

\begin{tabular}{ccl}
 &  & \tabularnewline
1.  & \includegraphics{actor}  & \multicolumn{1}{p{12cm}}{Presiona la opción, del menú, {[}Gestión de Proveedores{]}.}\tabularnewline
2.  & \includegraphics{sistema}  & \multicolumn{1}{p{12cm}}{Carga los datos, de la base de datos, de los proveedores existentes en el sistema {[}Trayectoria A{]}}\tabularnewline
3.  & \includegraphics{sistema}  & \multicolumn{1}{p{12cm}}{Despliega la pantalla PC1}\tabularnewline
4.  & \includegraphics{actor}  & \multicolumn{1}{p{12cm}}{Da clic en el botón {[}Registrar Proveedor{]} {[}Trayectoria B{]} {[}Trayectoria C{]} {[}Trayectoria D{]} {[}Trayectoria E{]}}\tabularnewline
5.  & \includegraphics{sistema}  & \multicolumn{1}{p{12cm}}{Extiende al CUC2.1 ``Registrar Proveedor"}\tabularnewline
 &  & \multicolumn{1}{p{12cm}}{.... Fin del caso de uso}\tabularnewline
\end{tabular}
\\


\textit{Trayectoria Alternativa A}

Condición: Ocurrió un error al conectar con la base de datos y no se pudieron cargar los proveedores registrados.

\begin{tabular}{ccl}
 &  & \tabularnewline
1.  & \includegraphics{sistema}  & \multicolumn{1}{p{12cm}}{ Muestra el mensaje de error MSG4 'No se pudo establecer conexión con la base de datos'}\tabularnewline
 &  & \multicolumn{1}{p{12cm}}{.... Fin de la trayectoria}\tabularnewline
\end{tabular}
\\

\textit{Trayectoria Alternativa B}

Condición: El actor escribe algún texto en el campo de búsqueda y presiona enter.

\begin{tabular}{ccl}
 &  & \tabularnewline
1.  & \includegraphics{sistema}  & \multicolumn{1}{p{12cm}}{Incluye al CUC2.0.1'}\tabularnewline
 &  & \multicolumn{1}{p{12cm}}{.... Fin de la trayectoria}\tabularnewline
\end{tabular}
\\


\textit{Trayectoria Alternativa C}

Condición: El actor pulsa el ícono de {[}Modificar Proveedor{]}

\begin{tabular}{ccl}
 &  & \tabularnewline
1.  & \includegraphics{sistema}  & \multicolumn{1}{p{12cm}}{ Extiende al CUC2.2: Modificar Proveedor.}\tabularnewline
 &  & \multicolumn{1}{p{12cm}}{.... Fin de la trayectoria}\tabularnewline
\end{tabular}
\\


\textit{Trayectoria Alternativa D}

Condición: El actor pulsa el ícono de {[}Eliminar Proveedor{]}

\begin{tabular}{ccl}
 &  & \tabularnewline
1.  & \includegraphics{sistema}  & \multicolumn{1}{p{12cm}}{ Extiende al CUC2.3: Eliminar Proveedor.}\tabularnewline
 &  & \multicolumn{1}{p{12cm}}{.... Fin de la trayectoria}\tabularnewline
\end{tabular}
\\


\textit{Trayectoria Alternativa E}

Condición: El actor pulsa el ícono de {[}Desbloquear Proveedor{]}

\begin{tabular}{ccl}
 &  & \tabularnewline
1.  & \includegraphics{sistema}  & \multicolumn{1}{p{12cm}}{ Extiende al CUC2.4: Desbloquear Proveedor.}\tabularnewline
 &  & \multicolumn{1}{p{12cm}}{.... Fin de la trayectoria}\tabularnewline
\end{tabular}

% = = = = = = = = = = = = = = = = = = = = = = = = =
%    FIN DE TRAYECTORIA
% = = = = = = = = = = = = = = = = = = = = = = = = = 

\newpage

% = = = = = = = = = = = = = = = = = = = = = = = = =
%    INICIO DE TRAYECTORIA
% = = = = = = = = = = = = = = = = = = = = = = = = =
\subsubsection{CUC2.0.1: Buscar}

\textbf{\large Resumen}\\
 {\large {} {} Este caso de uso tiene como objetivo permitir al
usuario de compras buscar un proveedor por medio palabra(s).}\\


% = = = = = = = = = = = = = = = = = = = = = = = = =
%    INICIO DE TABLA
% = = = = = = = = = = = = = = = = = = = = = = = = = 
\begin{table}[!ht]
\begin{centering}
\begin{tabular}{|c||c|l|}
\hline 
\multicolumn{2}{|c|}{Caso de Uso:} & CUC2.0.1: Buscar\tabularnewline
\hline 
\multicolumn{3}{|>{\columncolor[gray]{0.7}}c}{Resumen de Atributos}\tabularnewline
\hline 
\multicolumn{2}{|c|}{Autor} & \multicolumn{1}{p{10cm}||}{Jiménez López Carlos Adrian}\tabularnewline
\hline 
\multicolumn{2}{|c|}{Actor} & \multicolumn{1}{p{10cm}||}{Usuario de Compras}\tabularnewline
\hline 
\multicolumn{2}{|c|}{Propósito} & \multicolumn{1}{p{10cm}||}{Buscar los proveedores que coincidan con un criterio de búsqueda}\tabularnewline
\hline 
\multicolumn{2}{|c|}{Entradas} & \multicolumn{1}{p{10cm}||}{Palabra(s) Clave}\tabularnewline
\hline 
\multicolumn{2}{|c|}{Salidas} & \multicolumn{1}{p{10cm}||}{Lista de proveedores que coinciden con la entrada}\tabularnewline
\hline 
\multicolumn{2}{|c|}{Pre-condiciones} & \multicolumn{1}{p{10cm}||}{Iniciar sesión como usuario de compras}\tabularnewline
\hline 
\multicolumn{2}{|c|}{Pos-condiciones} & \multicolumn{1}{p{10cm}||}{Ninguna}\tabularnewline
\hline 
\multicolumn{2}{|c|}{Errores} & \multicolumn{1}{p{10cm}||}{MSG3 Datos no encontrados, MSG4 Error al conectar con la base de datos.}\tabularnewline
\hline 
\multicolumn{2}{|c|}{Tipo} & \multicolumn{1}{p{10cm}||}{Secundario}\tabularnewline
\hline 
\multicolumn{2}{|c|}{Fuente} & \multicolumn{1}{p{10cm}||}{Proveniente del CUC2.0.}\tabularnewline
\hline 
\end{tabular}
\par\end{centering}

\caption{Caso de Uso Compras 2.0.1: Buscar}


\label{CUC1.0} 
\end{table}


% = = = = = = = = = = = = = = = = = = = = = = = = =
%    FIN DE TABLA
% = = = = = = = = = = = = = = = = = = = = = = = = =


\textbf{\large Trayectorias del CU}{\large \par}

\textit{\large Trayectoria Principal}{\large {} }{\large \par}

\begin{tabular}{ccl}
 &  & \tabularnewline
1.  & \includegraphics{sistema}  & \multicolumn{1}{p{12cm}}{Verifica si la entrada no es vacía {[}Trayectoria
A{]}}\tabularnewline
2.  & \includegraphics{sistema}  & \multicolumn{1}{p{12cm}}{Consulta a los proveedores que cumplen con el criterio.{[}Trayectoria B{]}}\tabularnewline
3.  & \includegraphics{sistema}  & \multicolumn{1}{p{12cm}}{Muestra el resultado de la búsqueda{[}Trayectoria C{]}}\tabularnewline
 &  & \multicolumn{1}{p{12cm}}{.... Fin del caso de uso}\tabularnewline
\end{tabular}
\\

\textit{Trayectoria Alternativa A}

Condición: La entrada es vacía.

\begin{tabular}{ccl}
 &  & \tabularnewline
1.  & \includegraphics{sistema}  & \multicolumn{1}{p{12cm}}{Consulta todos los proveedores que existen.}\tabularnewline
2.  & \includegraphics{sistema}  & \multicolumn{1}{p{12cm}}{Muestra el resultado.}\tabularnewline
 &  & \multicolumn{1}{p{12cm}}{.... Fin de la trayectoria}\tabularnewline
\end{tabular}

\textit{Trayectoria Alternativa B}

Condición: Ocurrió un error al conectar con la base de datos y no se pudieron cargar los proveedores registrados.

\begin{tabular}{ccl}
 &  & \tabularnewline
1.  & \includegraphics{sistema}  & \multicolumn{1}{p{12cm}}{ Muestra el mensaje de error MSG4 'No se pudo establecer conexión
con la base de datos'}\tabularnewline
 &  & \multicolumn{1}{p{12cm}}{.... Fin de la trayectoria}\tabularnewline
\end{tabular}
\\

\textit{Trayectoria Alternativa C}

Condición: No arrojó resultados la búsqueda

\begin{tabular}{ccl}
 &  & \tabularnewline
1.  & \includegraphics{sistema}  & \multicolumn{1}{p{12cm}}{ Muestra el mensaje de error MSG3 'La búsqueda XXXXX no muestra ningún resultado'}\tabularnewline
 &  & \multicolumn{1}{p{12cm}}{.... Fin de la trayectoria}\tabularnewline
\end{tabular}
\\



% = = = = = = = = = = = = = = = = = = = = = = = = =
%    FIN DE TRAYECTORIA
% = = = = = = = = = = = = = = = = = = = = = = = = = 

\newpage

% = = = = = = = = = = = = = = = = = = = = = = = = =
%    INICIO DE TRAYECTORIA
% = = = = = = = = = = = = = = = = = = = = = = = = =



\subsubsection{CUC2.1: Registrar Proveedor}

\textbf{\large Resumen}\\
 {\large {} {} Este caso de uso tiene como objetivo permitir al
usuario de compras registrar un nuevo proveedor dentro del sistema.}\\


% = = = = = = = = = = = = = = = = = = = = = = = = =
%	INICIO DE TABLA
% = = = = = = = = = = = = = = = = = = = = = = = = = 
% Se utiliza para que no se salga la tabla de la hoja, porque sino se autodimensiona
%\multicolumn{1}{p{10cm}||}{Contenido}
\begin{table}[!ht]
\begin{centering}
\begin{tabular}{|c||c|l|}
\hline 
\multicolumn{2}{|c|}{Caso de Uso:} & CUC2.1: Registrar Proveedor\tabularnewline
\hline 
\multicolumn{3}{|>{\columncolor[gray]{0.7}}c}{Resumen de Atributos}\tabularnewline
\hline 
\multicolumn{2}{|c|}{Autor} & \multicolumn{1}{p{10cm}||}{Jiménez López Carlos Adrian}\tabularnewline
\hline 
\multicolumn{2}{|c|}{Actor} & \multicolumn{1}{p{10cm}||}{Usuario de Compras}\tabularnewline
\hline 
\multicolumn{2}{|c|}{Propósito} & \multicolumn{1}{p{10cm}||}{Registrar un nuevo proveedor }\tabularnewline
\hline 
\multicolumn{2}{|c|}{Entradas} & \multicolumn{1}{p{10cm}||}{RFC,
 Telefono, E-mail, Lista de productos.}\tabularnewline
\hline 
\multicolumn{2}{|c|}{Salidas} & \multicolumn{1}{p{10cm}||}{MSG1 Mensaje de confirmación.}\tabularnewline
\hline 
\multicolumn{2}{|c|}{Pre-condiciones} & \multicolumn{1}{p{10cm}||}{Ninguna}\tabularnewline
\hline 
\multicolumn{2}{|c|}{Pos-condiciones} & \multicolumn{1}{p{10cm}||}{El proveedor debe ser agregado a la base de datos}\tabularnewline
\hline 
\multicolumn{2}{|c|}{Errores} & \multicolumn{1}{p{10cm}||}{MSG4 Error al conectar con la base de datos. MSG7 Formato incorrecto
de los datos. MSG8 Campos obligatorios. MSG11 Datos duplicados.}\tabularnewline
\hline 
\multicolumn{2}{|c|}{Tipo} & \multicolumn{1}{p{10cm}||}{Secundario}\tabularnewline
\hline 
\multicolumn{2}{|c|}{Fuente} & \multicolumn{1}{p{10cm}||}{Proveniente del CUC2.0.}\tabularnewline
\hline 
\end{tabular}
\par\end{centering}

\caption{Caso de Uso Compras 2.1: Registrar Proveedor}


\label{CUC1.1} 
\end{table}


% = = = = = = = = = = = = = = = = = = = = = = = = =
%	FIN DE TABLA
% = = = = = = = = = = = = = = = = = = = = = = = = =


\textbf{\large Trayectorias del CU}{\large \par}

\textit{\large Trayectoria Principal}{\large {} }{\large \par}

%Solo hay que cambiar el nombre de la imagen dependiendo de si es actor o sistema


\begin{tabular}{ccl}
 &  & \tabularnewline
1.  & \includegraphics{actor}  & \multicolumn{1}{p{12cm}}{ Presiona el botón {[}Registrar Proveedor{]} en la pantalla PC1}\tabularnewline
2.  & \includegraphics{sistema}  & \multicolumn{1}{p{12cm}}{ Despliega la pantalla PC5 y pide los datos necesarios.}\tabularnewline
3.  & \includegraphics{actor}  & \multicolumn{1}{p{12cm}}{ Llena el formulario y presiona el botón {[}Registrar{]}}.\tabularnewline
4.  & \includegraphics{sistema}  & \multicolumn{1}{p{12cm}}{ Verifica que no haya campos vacíos {[}Trayectoria A{]}}\tabularnewline
5.  & \includegraphics{sistema}  & \multicolumn{1}{p{12cm}}{ Verifica que los campos se hayan llenado correctamente {[}Trayectoria B{]}}\tabularnewline
6.  & \includegraphics{sistema}  & \multicolumn{1}{p{12cm}}{ Verifica que el RFC ingresado no se encuentre ya registrado {[}Trayectoria C{]}}\tabularnewline
7.  & \includegraphics{sistema}  & \multicolumn{1}{p{12cm}}{ Registra al nuevo proveedor {[}Trayectoria D{]}}\tabularnewline
8.  & \includegraphics{sistema}  & \multicolumn{1}{p{12cm}}{ Despliega el mensaje MSG 1 'El proveedor ha sido registrado exitosamente'}\tabularnewline
9.  & \includegraphics{sistema}  & \multicolumn{1}{p{12cm}}{ Despliega la pantalla PC1}\tabularnewline
 &  & \multicolumn{1}{p{12cm}}{.... Fin del caso de uso}\tabularnewline
\end{tabular}\\

\textit{Trayectoria A}

Condición: Alguno de los campos está vacío y es obligatorio

\begin{tabular}{ccl}
 &  & \tabularnewline
1.  & \includegraphics{sistema}  & \multicolumn{1}{p{12cm}}{ Muestra el mensaje de error MSG8 'El campo RFC es un campo obligatorio',
'El campo teléfono es un campo obligatorio', 'El campo lista de productos
es un campo obligatorio' según corresponda}\tabularnewline
2.  & \includegraphics{sistema}  & \multicolumn{1}{p{12cm}}{Regresa a la pantalla PC5}\tabularnewline
 &  & \multicolumn{1}{p{12cm}}{.... Fin de la trayectoria}\tabularnewline
\end{tabular}\\


\textit{Trayectoria B}

Condición: Alguno de los campos fue llenado incorrectamente

\begin{tabular}{ccl}
 &  & \tabularnewline
1.  & \includegraphics{sistema}  & \multicolumn{1}{p{12cm}}{ Muestra el mensaje de error MSG7 'El RFC sólo acepta letras y números',
'El teléfono sólo acepta números', 'El email sólo acepta letras, numeros,
punto(.), guiones(-) y guión bajo(\_)' según corresponda}\tabularnewline
2.  & \includegraphics{sistema}  & \multicolumn{1}{p{12cm}}{Regresa a la pantalla PC5}\tabularnewline
 &  & \multicolumn{1}{p{12cm}}{.... Fin de la trayectoria}\tabularnewline
\end{tabular}\\

\textit{Trayectoria C}

Condición: El RFC ya se encuentra registrado

\begin{tabular}{ccl}
 &  & \tabularnewline
1.  & \includegraphics{sistema}  & \multicolumn{1}{p{12cm}}{ Muestra el mensaje de error MSG11 'El RFC ABCD910123 ya existe'}\tabularnewline
2.  & \includegraphics{sistema}  & \multicolumn{1}{p{12cm}}{Regresa a la pantalla PC5}\tabularnewline
 &  & \multicolumn{1}{p{12cm}}{.... Fin de la trayectoria}\tabularnewline
\end{tabular}\\


\textit{Trayectoria D}

Condición: Ocurrió un error al conectar con la base de datos y no se pudo registrar el proveedor en la base de datos

\begin{tabular}{ccl}
 &  & \tabularnewline
1.  & \includegraphics{sistema}  & \multicolumn{1}{p{12cm}}{Muestra el mensaje de error MSG4 'No se pudo establecer conexion con
la base de datos'}\tabularnewline
2.  & \includegraphics{sistema}  & \multicolumn{1}{p{12cm}}{Regresa a la pantalla PC5}\tabularnewline
 &  & \multicolumn{1}{p{12cm}}{.... Fin de la trayectoria}\tabularnewline
\end{tabular}

% = = = = = = = = = = = = = = = = = = = = = = = = =
%    FIN DE TRAYECTORIA
% = = = = = = = = = = = = = = = = = = = = = = = = = 

\newpage

% = = = = = = = = = = = = = = = = = = = = = = = = =
%    INICIO DE TRAYECTORIA
% = = = = = = = = = = = = = = = = = = = = = = = = =


\subsubsection{CUC2.2: Modificar Proveedor}

\textbf{\large Resumen}\\
 {\large {} {} Este caso de uso tiene como objetivo permitir al
usuario de compras modificar los datos de un proveedor.}\\


% = = = = = = = = = = = = = = = = = = = = = = = = =
%	INICIO DE TABLA
% = = = = = = = = = = = = = = = = = = = = = = = = = 
% Se utiliza para que no se salga la tabla de la hoja, porque sino se autodimensiona
%\multicolumn{1}{p{10cm}||}{Contenido}
\begin{table}[!ht]
\begin{centering}
\begin{tabular}{|c||c|l|}
\hline 
\multicolumn{2}{|c|}{Caso de Uso:} & CUC2.2: Modificar Proveedor\tabularnewline
\hline 
\multicolumn{3}{|>{\columncolor[gray]{0.7}}c}{Resumen de Atributos}\tabularnewline
\hline 
\multicolumn{2}{|c|}{Autor} & \multicolumn{1}{p{10cm}||}{Jiménez López Carlos Adrian}\tabularnewline
\hline 
\multicolumn{2}{|c|}{Actor} & \multicolumn{1}{p{10cm}||}{Usuario de Compras}\tabularnewline
\hline 
\multicolumn{2}{|c|}{Propósito} & \multicolumn{1}{p{10cm}||}{Modificar los datos de un proveedor registrado}\tabularnewline
\hline 
\multicolumn{2}{|c|}{Entradas} & \multicolumn{1}{p{10cm}||}{RFC, 
 Telefono, E-mail, Lista de productos}\tabularnewline
\hline 
\multicolumn{2}{|c|}{Salidas} & \multicolumn{1}{p{10cm}||}{MSG1 Mensaje de confirmación.}\tabularnewline
\hline 
\multicolumn{2}{|c|}{Pre-condiciones} & \multicolumn{1}{p{10cm}||}{Debe seleccionar un proveedor de la pantalla PC1}\tabularnewline
\hline 
\multicolumn{2}{|c|}{Pos-condiciones} & \multicolumn{1}{p{10cm}||}{El proveedor debe ser actualizado en la base de datos}\tabularnewline
\hline 
\multicolumn{2}{|c|}{Errores} & \multicolumn{1}{p{10cm}||}{MSG4 Error al conectar con la base de datos. MSG7 Formato incorrecto
de los datos. MSG8 Campos obligatorios. MSG11 Datos duplicados.}\tabularnewline
\hline 
\multicolumn{2}{|c|}{Tipo} & \multicolumn{1}{p{10cm}||}{Secundario}\tabularnewline
\hline 
\multicolumn{2}{|c|}{Fuente} & \multicolumn{1}{p{10cm}||}{Proveniente del CUC2.0.}\tabularnewline
\hline 
\end{tabular}
\par\end{centering}

\caption{Caso de Uso Compras 2.2: Modificar Proveedor}


\label{CUC1.2} 
\end{table}


% = = = = = = = = = = = = = = = = = = = = = = = = =
%	FIN DE TABLA
% = = = = = = = = = = = = = = = = = = = = = = = = =


\textbf{\large Trayectorias del CU}{\large \par}

\textit{\large Trayectoria Principal}{\large {} }{\large \par}

%Solo hay que cambiar el nombre de la imagen dependiendo de si es actor o sistema


\begin{tabular}{ccl}
 &  & \tabularnewline
1.  & \includegraphics{actor}  & \multicolumn{1}{p{12cm}}{ Presiona el ícono {[}Modificar Proveedor{]} en la pantalla PC1}\tabularnewline
2.  & \includegraphics{sistema}  & \multicolumn{1}{p{12cm}}{ Despliega la pantalla PC6 y pide los datos necesarios.}\tabularnewline
3.  & \includegraphics{actor}  & \multicolumn{1}{p{12cm}}{ Cambia los datos que desea modificar y presiona el botón {[}Actualizar{]}}.\tabularnewline
4.  & \includegraphics{sistema}  & \multicolumn{1}{p{12cm}}{ Verifica que no haya campos vacíos {[}Trayectoria A{]}}\tabularnewline
5.  & \includegraphics{sistema}  & \multicolumn{1}{p{12cm}}{ Verifica que los campos se hayan llenado correctamente {[}Trayectoria B{]}}\tabularnewline
6.  & \includegraphics{sistema}  & \multicolumn{1}{p{12cm}}{ Verifica que el RFC ingresado no se encuentre ya registrado {[}Trayectoria C{]}}\tabularnewline
7.  & \includegraphics{sistema}  & \multicolumn{1}{p{12cm}}{ Actualiza los datos del proveedor {[}Trayectoria D{]}}\tabularnewline
8.  & \includegraphics{sistema}  & \multicolumn{1}{p{12cm}}{ Despliega el mensaje MSG 1 'El proveedor ha sido modificado exitosamente'}\tabularnewline
9.  & \includegraphics{sistema}  & \multicolumn{1}{p{12cm}}{ Despliega la pantalla PC1}\tabularnewline
 &  & \multicolumn{1}{p{12cm}}{.... Fin del caso de uso}\tabularnewline
\end{tabular}\\


\textit{Trayectoria A}

Condición: Alguno de los campos está vacío y es obligatorio

\begin{tabular}{ccl}
 &  & \tabularnewline
1.  & \includegraphics{sistema}  & \multicolumn{1}{p{12cm}}{ Muestra el mensaje de error MSG8 'El campo RFC es un campo obligatorio',
'El campo teléfono es un campo obligatorio', 'El campo lista de productos
es un campo obligatorio' según corresponda}\tabularnewline
2.  & \includegraphics{sistema}  & \multicolumn{1}{p{12cm}}{Regresa a la pantalla PC6}\tabularnewline
 &  & \multicolumn{1}{p{12cm}}{.... Fin de la trayectoria}\tabularnewline
\end{tabular}\\


\textit{Trayectoria B}

Condición: Alguno de los campos fue llenado incorrectamente

\begin{tabular}{ccl}
 &  & \tabularnewline
1.  & \includegraphics{sistema}  & \multicolumn{1}{p{12cm}}{ Muestra el mensaje de error MSG7 'El RFC sólo acepta letras y números',
'El teléfono sólo acepta números', 'El email sólo acepta letras, numeros,
punto(.), guiones(-) y guión bajo(\_)' según corresponda}\tabularnewline
2.  & \includegraphics{sistema}  & \multicolumn{1}{p{12cm}}{Regresa a la pantalla PC6}\tabularnewline
 &  & \multicolumn{1}{p{12cm}}{.... Fin de la trayectoria}\tabularnewline
\end{tabular}\\

\textit{Trayectoria C}

Condición: El RFC ya se encuentra registrado

\begin{tabular}{ccl}
 &  & \tabularnewline
1.  & \includegraphics{sistema}  & \multicolumn{1}{p{12cm}}{ Muestra el mensaje de error MSG11 'El RFC ABCD910123 ya existe'}\tabularnewline
2.  & \includegraphics{sistema}  & \multicolumn{1}{p{12cm}}{Regresa a la pantalla PC6}\tabularnewline
 &  & \multicolumn{1}{p{12cm}}{.... Fin de la trayectoria}\tabularnewline
\end{tabular}\\


\textit{Trayectoria D}

Condición: Ocurrió un error al conectar con la base de datos y no se pudo actualizar el proveedor en la base de datos

\begin{tabular}{ccl}
 &  & \tabularnewline
1.  & \includegraphics{sistema}  & \multicolumn{1}{p{12cm}}{Muestra el mensaje de error MSG4 'No se pudo establecer conexion con
la base de datos'}\tabularnewline
2.  & \includegraphics{sistema}  & \multicolumn{1}{p{12cm}}{Regresa a la pantalla PC6}\tabularnewline
 &  & \multicolumn{1}{p{12cm}}{.... Fin de la trayectoria}\tabularnewline
\end{tabular}

% = = = = = = = = = = = = = = = = = = = = = = = = =
%    FIN DE TRAYECTORIA
% = = = = = = = = = = = = = = = = = = = = = = = = = 

\newpage

% = = = = = = = = = = = = = = = = = = = = = = = = =
%    INICIO DE TRAYECTORIA
% = = = = = = = = = = = = = = = = = = = = = = = = =


\subsubsection{CUC2.3: Eliminar Proveedor}

\textbf{\large Resumen}\\
 {\large {} {} Este caso de uso tiene como objetivo permitir al
usuario de compras eliminar los datos de un proveedor.}\\


% = = = = = = = = = = = = = = = = = = = = = = = = =
%	INICIO DE TABLA
% = = = = = = = = = = = = = = = = = = = = = = = = = 
% Se utiliza para que no se salga la tabla de la hoja, porque sino se autodimensiona
%\multicolumn{1}{p{10cm}||}{Contenido}
\begin{table}[!ht]
\begin{centering}
\begin{tabular}{|c||c|l|}
\hline 
\multicolumn{2}{|c|}{Caso de Uso:} & CUC2.3: Eliminar Proveedor\tabularnewline
\hline 
\multicolumn{3}{|>{\columncolor[gray]{0.7}}c}{Resumen de Atributos}\tabularnewline
\hline 
\multicolumn{2}{|c|}{Autor} & \multicolumn{1}{p{10cm}||}{Jiménez López Carlos Adrian}\tabularnewline
\hline 
\multicolumn{2}{|c|}{Actor} & \multicolumn{1}{p{10cm}||}{Usuario de Compras}\tabularnewline
\hline 
\multicolumn{2}{|c|}{Propósito} & \multicolumn{1}{p{10cm}||}{Eliminar un proveedor de manera permanente}\tabularnewline
\hline 
\multicolumn{2}{|c|}{Entradas} & \multicolumn{1}{p{10cm}||}{Ninguna}\tabularnewline
\hline 
\multicolumn{2}{|c|}{Salidas} & \multicolumn{1}{p{10cm}||}{MSG1 Mensaje de confirmación.}\tabularnewline
\hline 
\multicolumn{2}{|c|}{Pre-condiciones} & \multicolumn{1}{p{10cm}||}{Debe seleccionar un proveedor de la pantalla PC1}\tabularnewline
\hline 
\multicolumn{2}{|c|}{Pos-condiciones} & \multicolumn{1}{p{10cm}||}{El proveedor debe ser eliminado de la base de datos}\tabularnewline
\hline 
\multicolumn{2}{|c|}{Errores} & \multicolumn{1}{p{10cm}||}{MSG4 Error al conectar con la base de datos.}\tabularnewline
\hline 
\multicolumn{2}{|c|}{Tipo} & \multicolumn{1}{p{10cm}||}{Secundario}\tabularnewline
\hline 
\multicolumn{2}{|c|}{Fuente} & \multicolumn{1}{p{10cm}||}{Proveniente del CUC2.0.}\tabularnewline
\hline 
\end{tabular}
\par\end{centering}

\caption{Caso de Uso Compras 2.3: Eliminar Proveedor}


\label{CUC1.3} 
\end{table}


% = = = = = = = = = = = = = = = = = = = = = = = = =
%	FIN DE TABLA
% = = = = = = = = = = = = = = = = = = = = = = = = =


\textbf{\large Trayectorias del CU}{\large \par}

\textit{\large Trayectoria Principal}{\large {} }{\large \par}

%Solo hay que cambiar el nombre de la imagen dependiendo de si es actor o sistema


\begin{tabular}{ccl}
 &  & \tabularnewline
1.  & \includegraphics{actor}  & \multicolumn{1}{p{12cm}}{ Presiona el ícono {[}Eliminar Proveedor{]} en la pantalla PC1}\tabularnewline
2.  & \includegraphics{sistema}  & \multicolumn{1}{p{12cm}}{Despliega el MSG9 '¿Desea eliminar el proveedor?' pidiendo la confirmación
para eliminar.}\tabularnewline
3.  & \includegraphics{actor}  & \multicolumn{1}{p{12cm}}{Presiona el botón {[}Eliminar{]} {[}Trayectoria A{]}}\tabularnewline
4.  & \includegraphics{sistema}  & \multicolumn{1}{p{12cm}}{ Elimina al proveedor {[}Trayectoria B{]}}\tabularnewline
5.  & \includegraphics{sistema}  & \multicolumn{1}{p{12cm}}{ Despliega el mensaje MSG 1 'El proveedor ha sido eliminado exitosamente'}\tabularnewline
6.  & \includegraphics{sistema}  & \multicolumn{1}{p{12cm}}{ Despliega la pantalla PC1}\tabularnewline
 &  & \multicolumn{1}{p{12cm}}{.... Fin del caso de uso}\tabularnewline
\end{tabular}\\


\textit{Trayectoria A}

Condición: El usuario presiona el botón {[}Cancelar{]}

\begin{tabular}{ccl}
 &  & \tabularnewline
1.  & \includegraphics{sistema}  & \multicolumn{1}{p{12cm}}{Regresa a la pantalla PC1}\tabularnewline
 &  & \multicolumn{1}{p{12cm}}{.... Fin de la trayectoria}\tabularnewline
\end{tabular}\\


\textit{Trayectoria B}

Condición: Ocurrió un error al conectar con la base de datos y no se pudo eliminar el proveedor de la base de datos

\begin{tabular}{ccl}
 &  & \tabularnewline
1.  & \includegraphics{sistema}  & \multicolumn{1}{p{12cm}}{Muestra el mensaje de error MSG4 'No se pudo establecer conexion con
la base de datos'}\tabularnewline
2.  & \includegraphics{sistema}  & \multicolumn{1}{p{12cm}}{Regresa a la pantalla PC1}\tabularnewline
 &  & \multicolumn{1}{p{12cm}}{.... Fin de la trayectoria}\tabularnewline
\end{tabular}

% = = = = = = = = = = = = = = = = = = = = = = = = =
%    FIN DE TRAYECTORIA
% = = = = = = = = = = = = = = = = = = = = = = = = = 

\newpage

% = = = = = = = = = = = = = = = = = = = = = = = = =
%    INICIO DE TRAYECTORIA
% = = = = = = = = = = = = = = = = = = = = = = = = =


\subsubsection{CUC2.4: Desbloquear Proveedor}

\textbf{\large Resumen}\\
 {\large {} {} Este caso de uso tiene como objetivo permitir al
usuario de compras desbloquear a un proveedor. Ver BR}\\


% = = = = = = = = = = = = = = = = = = = = = = = = =
%	INICIO DE TABLA
% = = = = = = = = = = = = = = = = = = = = = = = = = 
% Se utiliza para que no se salga la tabla de la hoja, porque sino se autodimensiona
%\multicolumn{1}{p{10cm}||}{Contenido}
\begin{table}[!ht]
\begin{centering}
\begin{tabular}{|c||c|l|}
\hline 
\multicolumn{2}{|c|}{Caso de Uso:} & CUC2.4: Desbloquear Proveedor\tabularnewline
\hline 
\multicolumn{3}{|>{\columncolor[gray]{0.7}}c}{Resumen de Atributos}\tabularnewline
\hline 
\multicolumn{2}{|c|}{Autor} & \multicolumn{1}{p{10cm}||}{Jiménez López Carlos Adrian}\tabularnewline
\hline 
\multicolumn{2}{|c|}{Actor} & \multicolumn{1}{p{10cm}||}{Usuario de Compras}\tabularnewline
\hline 
\multicolumn{2}{|c|}{Propósito} & \multicolumn{1}{p{10cm}||}{Desbloquear un proveedor}\tabularnewline
\hline 
\multicolumn{2}{|c|}{Entradas} & \multicolumn{1}{p{10cm}||}{Ninguna}\tabularnewline
\hline 
\multicolumn{2}{|c|}{Salidas} & \multicolumn{1}{p{10cm}||}{MSG1 Mensaje de confirmación.}\tabularnewline
\hline 
\multicolumn{2}{|c|}{Pre-condiciones} & \multicolumn{1}{p{10cm}||}{Debe seleccionar un proveedor de la pantalla PC1.\medskip{}
 El proveedor debe estar bloqueado}\tabularnewline
\hline 
\multicolumn{2}{|c|}{Pos-condiciones} & \multicolumn{1}{p{10cm}||}{El proveedor debe ser desbloqueado de la base de datos}\tabularnewline
\hline 
\multicolumn{2}{|c|}{Errores} & \multicolumn{1}{p{10cm}||}{MSG4 Error al conectar con la base de datos.}\tabularnewline
\hline 
\multicolumn{2}{|c|}{Tipo} & \multicolumn{1}{p{10cm}||}{Secundario}\tabularnewline
\hline 
\multicolumn{2}{|c|}{Fuente} & \multicolumn{1}{p{10cm}||}{Proveniente del CUC2.0.}\tabularnewline
\hline 
\end{tabular}
\par\end{centering}

\caption{Caso de Uso Compras 2.4: Desbloquear Proveedor}


\label{tab:CasosdeUso:nombredecasodeuso-1-1-1} 
\end{table}


% = = = = = = = = = = = = = = = = = = = = = = = = =
%	FIN DE TABLA
% = = = = = = = = = = = = = = = = = = = = = = = = =


\textbf{\large Trayectorias del CU}{\large \par}

\textit{\large Trayectoria Principal}{\large {} }{\large \par}

%Solo hay que cambiar el nombre de la imagen dependiendo de si es actor o sistema


\begin{tabular}{ccl}
 &  & \tabularnewline
1.  & \includegraphics{actor}  & \multicolumn{1}{p{12cm}}{ Presiona el ícono {[}Desbloquear Proveedor{]} en la pantalla PC1}\tabularnewline
2.  & \includegraphics{sistema}  & \multicolumn{1}{p{12cm}}{Despliega el MSG9 '¿Desea desbloquear el proveedor?' pidiendo la confirmación
para desbloquear.}\tabularnewline
3.  & \includegraphics{actor}  & \multicolumn{1}{p{12cm}}{Presiona el botón {[}Desbloquear{]} {[}Trayectoria A{]}}\tabularnewline
4.  & \includegraphics{sistema}  & \multicolumn{1}{p{12cm}}{ Desbloquea al proveedor {[}Trayectoria B{]}}\tabularnewline
5.  & \includegraphics{sistema}  & \multicolumn{1}{p{12cm}}{ Despliega el mensaje MSG 1 'El proveedor ha sido desbloqueado exitosamente'}\tabularnewline
6.  & \includegraphics{sistema}  & \multicolumn{1}{p{12cm}}{ Despliega la pantalla PC1}\tabularnewline
 &  & \multicolumn{1}{p{12cm}}{.... Fin del caso de uso}\tabularnewline
\end{tabular}\\


\textit{Trayectoria A}

Condición: El usuario presiona el botón {[}Cancelar{]}

\begin{tabular}{ccl}
 &  & \tabularnewline
1.  & \includegraphics{sistema}  & \multicolumn{1}{p{12cm}}{Regresa a la pantalla PC1}\tabularnewline
 &  & \multicolumn{1}{p{12cm}}{.... Fin de la trayectoria}\tabularnewline
\end{tabular}\\


\textit{Trayectoria B}

Condición: Ocurrió un error al conectar con la base de datos y no se pudo desbloquear el proveedor en la base de datos

\begin{tabular}{ccl}
 &  & \tabularnewline
1.  & \includegraphics{sistema}  & \multicolumn{1}{p{12cm}}{Muestra el mensaje de error MSG4 'No se pudo establecer conexion con
la base de datos'}\tabularnewline
2.  & \includegraphics{sistema}  & \multicolumn{1}{p{12cm}}{Regresa a la pantalla PC1}\tabularnewline
 &  & \multicolumn{1}{p{12cm}}{.... Fin de la trayectoria}\tabularnewline
\end{tabular}


\end{document}
